% GENERATED FILE. Edit canonical lessons, then run npm run generate:books.
% canonical-source-hash: fnv1a64:5ba53b7b074372d9
% canonical-lessons: GU-C273-arune-pidhi, GU-C273-gayo-che, GU-R273-who-did-it

\chapter{Who did it, and what has happened}
\label{ch:gu-who-did-it-and-what-has-happened}
\hlchaptermodality{273}

\begin{chapteropening}
\textbf{By the end of this chapter:} \emph{I can say who did something with a name or a noun (Aruṇe cā pīdhī, Mīrāe patra lakhyo), and say what has happened and what had happened (gayo che, karyũ hatũ).}

Complete GU-R273-who-did-it: say who did something, and what has and had happened.
\end{chapteropening}

\section[\texorpdfstring{\gu{અરુણે} \gu{ચા} \gu{પીધી} (Aruṇe cā pīdhī)}{અરુણે ચા પીધી (Aruṇe cā pīdhī)}]{\gu{અરુણે} \gu{ચા} \gu{પીધી} (Aruṇe cā pīdhī) — a name as the doer}

\label{lesson:GU-C273-arune-pidhi}



\begin{quote}
Before the new one: say \textquotedblleft{}I drank tea\textquotedblright{} and \textquotedblleft{}he drank water\textquotedblright{}. (\emph{mẽ cā pīdhī}, \emph{teṇe pāṇī pīdhũ}.) Both times the doer was a pronoun. This lesson puts a name in its place.
\end{quote}

\begin{grammarlens}[title={Grammar Lens: the doer takes -\gu{એ}}]
\begin{quote}
\textbf{\gu{અરુણે} \gu{ચા} \gu{પીધી}.}
\end{quote}

\begin{quote}
\emph{Aruṇe cā pīdhī.} — \textquotedblleft{}Arun drank tea.\textquotedblright{}
\end{quote}

With a verb that has an object, a name or a noun that does it takes \textbf{-\gu{એ}} (\emph{-e}), the way \textbf{\gu{હું}} becomes \textbf{\gu{મેં}}:

\noindent
\begin{tabularx}{\linewidth}{@{}>{\raggedright\arraybackslash}X>{\raggedright\arraybackslash}X@{}}
\toprule
\textbf{the doer} & \textbf{with -\gu{એ}} \\
\midrule
\textbf{\gu{અરુણ}} & \textbf{\gu{અરુણે}} (\emph{Aruṇe}) \\
\textbf{\gu{મીરા}} & \textbf{\gu{મીરાએ}} (\emph{Mīrāe}) \\
\textbf{\gu{બહેન}} & \textbf{\gu{બહેને}} (\emph{bahene}) \\
\textbf{\gu{પોલીસ}} & \textbf{\gu{પોલીસે}} (\emph{polīse}) \\
\textbf{\gu{છોકરો}} & \textbf{\gu{છોકરાએ}} (\emph{chokrāe}) \\
\bottomrule
\end{tabularx}

After a consonant it is the sign \textbf{\gu{ે}}; after a vowel, the letter \textbf{\gu{એ}}. \textbf{\gu{છોકરો}} first turns its \textbf{-\gu{ો}} into \textbf{-\gu{ા}}, as it does before \textbf{-\gu{ને}}. A plural keeps its \textbf{-\gu{ઓ}} and adds \textbf{\gu{એ}}: \textbf{\gu{છોકરાઓએ}} (\emph{chokrāoe}), \textquotedblleft{}the boys\textquotedblright{}.
\end{grammarlens}

\subsection*{You'll want to know: the verb follows the thing}
\begin{quote}
\textbf{\gu{મીરાએ} \gu{પત્ર} \gu{લખ્યો}.}
\end{quote}

\begin{quote}
\emph{Mīrāe patra lakhyo.} — \textquotedblleft{}Mira wrote a letter.\textquotedblright{}
\end{quote}

\textbf{\gu{પત્ર}} is masculine, so \textbf{\gu{લખ્યો}}. The doer leaves the ending alone; the thing sets it: \textbf{\gu{મીરાએ} \gu{વાર્તા} \gu{લખી}} (\emph{Mīrāe vārtā lakhī}), for a feminine story, and \textbf{\gu{અરુણે} \gu{પુસ્તક} \gu{વાંચ્યું}} (\emph{Aruṇe pustak vā̃cyũ}), for a neuter book.

When the thing is a person, it takes \textbf{-\gu{ને}}, and in Gujarati the verb follows that person too:

\begin{quote}
\textbf{\gu{અરુણે} \gu{મીરાને} \gu{જોઈ}.}
\end{quote}

\begin{quote}
\emph{Aruṇe Mīrāne joī.} — \textquotedblleft{}Arun saw Mira.\textquotedblright{}
\end{quote}

\textbf{\gu{જોઈ}} is feminine, for Mira. Turn it round, \textbf{\gu{મીરાએ} \gu{અરુણને} \gu{જોયો}} (\emph{Mīrāe Aruṇne joyo}): masculine, for Arun. Some speakers say \textbf{\gu{જોયું}} in both; this book follows the person.

\subsection*{Your turn}
Say these aloud:

\begin{itemize}
\raggedright
  \item \textquotedblleft{}The boy ate rice\textquotedblright{} — \emph{chokrāe bhāt khādho}
  \item \textquotedblleft{}My sister wrote a story\textquotedblright{} — \emph{mārī bahene vārtā lakhī}
  \item \textquotedblleft{}Mira saw Arun\textquotedblright{} — \emph{Mīrāe Aruṇne joyo}
\end{itemize}

\begin{tcolorbox}[breakable,colback=teal!4,colframe=teal!35!black,title={Before you move on}]
\textquotedblleft{}Arun drank tea\textquotedblright{}? (\textbf{\gu{અરુણે} \gu{ચા} \gu{પીધી}}.) Why \textbf{\gu{પીધી}}? (\textbf{\gu{ચા}} is feminine.) \textquotedblleft{}Mira wrote a letter\textquotedblright{}? (\textbf{\gu{મીરાએ} \gu{પત્ર} \gu{લખ્યો}}.) Who sets the ending, Mira or the letter? (The letter.)
\end{tcolorbox}
\section[\texorpdfstring{\gu{ગયો} \gu{છે} · \gu{ગયો} \gu{હતો} (gayo che · gayo hato)}{ગયો છે · ગયો હતો (gayo che · gayo hato)}]{\gu{ગયો} \gu{છે}, \gu{ગયો} \gu{હતો} (gayo che, gayo hato) — has gone, had gone}

\label{lesson:GU-C273-gayo-che}



\begin{quote}
Say \textquotedblleft{}he went\textquotedblright{} and \textquotedblleft{}she came\textquotedblright{}. (\emph{te gayo}, \emph{te āvī}.) Then \textquotedblleft{}he was\textquotedblright{}. (\emph{te hato}.) Put a past and a copula side by side, and a new past appears.
\end{quote}

\begin{grammarlens}[title={Grammar Lens: has gone}]
\begin{quote}
\textbf{\gu{અરુણ} \gu{બજાર} \gu{ગયો} \gu{છે}.}
\end{quote}

\begin{quote}
\emph{Aruṇ bajār gayo che.} — \textquotedblleft{}Arun has gone to the market.\textquotedblright{}
\end{quote}

\textbf{\gu{ગયો}} says he went. \textbf{\gu{ગયો} \gu{છે}} says he went and is still away: the going matters now. The \textbf{\gu{છે}} follows the person, as in the present: \textbf{\gu{હું} \gu{હમણાં} \gu{આવ્યો} \gu{છું}} (\emph{hũ hamṇā̃ āvyo chũ}), \textquotedblleft{}I have come this moment\textquotedblright{}; \textbf{\gu{અમે} \gu{આવ્યા} \gu{છીએ}} (\emph{ame āvyā chīe}), \textquotedblleft{}we have come\textquotedblright{}.

A verb with an object keeps its doer form and follows the thing: \textbf{\gu{મેં} \gu{કામ} \gu{કર્યું} \gu{છે}} (\emph{mẽ kām karyũ che}), \textquotedblleft{}I have done the work\textquotedblright{}. There the copula is \textbf{\gu{છે}}, whoever the doer.
\end{grammarlens}

\subsection*{You'll want to know: had gone}
\begin{quote}
\textbf{\gu{હું} \gu{સ્ટેશન} \gu{પહોંચ્યો}, \gu{પણ} \gu{ટ્રેન} \gu{ઊપડી} \gu{ગઈ} \gu{હતી}.}
\end{quote}

\begin{quote}
\emph{hũ sṭeśan pahõcyo, paṇ ṭren ūpaḍī gaī hatī.} — \textquotedblleft{}I reached the station, but the train had left.\textquotedblright{}
\end{quote}

\textbf{\gu{હતી}} puts the leaving further back, before another past. It follows the way \textbf{\gu{હતો}}, \textbf{\gu{હતી}}, \textbf{\gu{હતું}}, \textbf{\gu{હતા}} do:

\noindent
\begin{tabularx}{\linewidth}{@{}>{\raggedright\arraybackslash}X>{\raggedright\arraybackslash}X>{\raggedright\arraybackslash}X@{}}
\toprule
\textbf{} & \textbf{has (now)} & \textbf{had (before)} \\
\midrule
he went & \textbf{\gu{તે} \gu{ગયો} \gu{છે}} & \textbf{\gu{તે} \gu{ગયો} \gu{હતો}} \\
she went & \textbf{\gu{તે} \gu{ગઈ} \gu{છે}} & \textbf{\gu{તે} \gu{ગઈ} \gu{હતી}} \\
I did the work & \textbf{\gu{મેં} \gu{કામ} \gu{કર્યું} \gu{છે}} & \textbf{\gu{મેં} \gu{કામ} \gu{કર્યું} \gu{હતું}} \\
Mira drank tea & \textbf{\gu{મીરાએ} \gu{ચા} \gu{પીધી} \gu{છે}} & \textbf{\gu{મીરાએ} \gu{ચા} \gu{પીધી} \gu{હતી}} \\
\bottomrule
\end{tabularx}

\subsection*{Your turn}
Say these aloud:

\begin{itemize}
\raggedright
  \item \textquotedblleft{}We have come\textquotedblright{} — \emph{ame āvyā chīe}
  \item \textquotedblleft{}Arun had drunk the tea\textquotedblright{} — \emph{Aruṇe cā pīdhī hatī}
\end{itemize}

\begin{tcolorbox}[breakable,colback=teal!4,colframe=teal!35!black,title={Before you move on}]
\textquotedblleft{}He has gone\textquotedblright{}? (\textbf{\gu{તે} \gu{ગયો} \gu{છે}}.) \textquotedblleft{}She had gone\textquotedblright{}? (\textbf{\gu{તે} \gu{ગઈ} \gu{હતી}}.) \textquotedblleft{}I have done the work\textquotedblright{}? (\textbf{\gu{મેં} \gu{કામ} \gu{કર્યું} \gu{છે}}.) In \textbf{\gu{મીરાએ} \gu{ચા} \gu{પીધી} \gu{હતી}}, why \textbf{\gu{હતી}}? (It follows \textbf{\gu{ચા}}.)
\end{tcolorbox}
\section[(review)]{Review — who did it, and what has happened}

\label{lesson:GU-R273-who-did-it}



\begin{quote}
Say \textquotedblleft{}I drank tea\textquotedblright{} and \textquotedblleft{}she went home\textquotedblright{}. (\emph{mẽ cā pīdhī}, \emph{te ghare gaī}.)
\end{quote}

\subsection*{Your turn}
Say these aloud:

\begin{itemize}
\raggedright
  \item Arun ate rice. Mira drank water.
  \item My sister wrote a letter, and Arun read it.
  \item Mira saw Arun at the market.
  \item Arun has gone home. I have done the work.
  \item I reached the station, but the train had left.
\end{itemize}

\begin{tcolorbox}[breakable,colback=teal!4,colframe=teal!35!black,title={Before you move on}]
In \textbf{\gu{મીરાએ} \gu{પત્ર} \gu{લખ્યો}}, which word sets \textbf{\gu{લખ્યો}}? (\textbf{\gu{પત્ર}}.) Now the other way round: put Arun in Mira's place. Does the verb change? (No: \textbf{\gu{અરુણે} \gu{પત્ર} \gu{લખ્યો}}.) Put \textbf{\gu{વાર્તા}} in place of \textbf{\gu{પત્ર}}? (Yes: \textbf{\gu{લખી}}.) And \textquotedblleft{}had gone\textquotedblright{}, for a woman? (\textbf{\gu{ગઈ} \gu{હતી}}.)
\end{tcolorbox}
